\documentstyle[%


righttag%


,11pt%


,amssymb%


,russian%               remove in Eng.


,izvru%                 Set izven% in Eng.


% ,izvru-ks             for brief communications


]{amsart}





%      Math definitions





\newcommand{\A}{{\cal A}}


\newcommand{\col}{\operatorname{col}}


\newcommand{\vraisup}{\operatorname{vrai~sup}}


\newcommand{\tts}[1]{}





%\hoffset=-15mm                  For Eng.





\setcounter{page}{53}


\god{1998}


\nomer{4~(431)} %                Remove in Eng.


%\nomer{Vol.\,42, No.\,2}        Set in Eng.


% \str{\pageref{begin}--\pageref{end}}    Set in Eng.


\udk{517.988}


\author{\it И.Н.\,Майборода, Л.А.\,Островецкий}


% NB    ^^ remove in Eng.


\title{О двустороннем методе решения нелинейных уравнений}





\date{}


% \pagestyle{myheadings}         Set in Eng.





\pagestyle{plain} %              Remove in Eng.


\begin{document}





% \thispagestyle{plain}          Set in Eng.





\maketitle


\label{begin}





При приближенном решении операторных уравнений


очень удобны двусторонние алгоритмы. Они определяют


два приближения, аппроксимирующие искомое решение снизу и сверху,


что позволяет свести вопрос об оценке погрешности к простому вычислению


нормы разности найденных приближений. Обычно в этих алгоритмах


предполагается некоторая монотонность.





Теория решений уравнений с монотонными операторами при различных их


свойствах и свойствах пространств с конусами, в которых они рассматриваются,


развита в работах М.А.\,Красно\-сельского, И.А.\,Бахтина, В.Я.\,Стеценко и


других авторов [1]--[3].


Теория монотонных операторов В.И.\,Опойцевым была перенесена на


гетерогенные и гетеротонные операторы [4], [5].


В работах Н.С.\,Курпеля и других математиков [6] (в этой монографии имеется


подробная библиография) исследовались двусторонние конструкции немонотонных


операторов с производными, обладающими теми или иными свойствами


монотонности.





В данной статье доказано, что некоторые результаты, характерные для


вышеуказанных операторов и полученные в работах [4], [5], распространяются на


операторы с гетерогенной производной в банаховых пространствах с различными


конусами.





Напомним некоторые используемые в данной работе определения.





Пусть $B$ --- вещественное банахово пространство с конусом $K$ и


$\Im=\{P_\alpha\mid\alpha\in\Omega\}$ ---


некоторое семейство расщепляющих операторов. Природа множества


$\Omega$ безразлична.


При любом $\alpha\in\Omega$ операторы $P_\alpha$ действуют из


$B$ в $B$. Пусть $Q_\alpha=I-P_\alpha$, где $I$ ---


тождественное преобразование. Если два элемента


$x,y\in B$ принадлежат некоторому


конусному отрезку


$\langle u;v\rangle=\{x|u\le x\le v\}$, то при любом $\alpha$ имеет место


$P_\alpha x+Q_\alpha y\in\langle u;v\rangle$.





Пусть $\Re=\{R_\alpha\mid\alpha\in\Omega\}$ --- некоторое семейство


операторов, расщепляющих любой


элемент $x\in B$ на компоненты $x_\alpha$ ($x_\alpha=R_\alpha x$).


При этом каждый


оператор $R_\alpha$ отображает


$B$ в некоторое множество $G_\alpha$, а семейство $\Re$ каждому


$x\in B$ сопоставляет множество


компонент $\{x_\alpha\}$, причем $\{x_\alpha\}\ne\{y_\alpha\}$,


если $x\ne y$. Операция


объединения множества


компонент $\{x_\alpha\}$ в элемент $x\in B$ обозначается


$\sum\limits_\alpha\{R_\alpha x\}=\sum\limits_\alpha\{x_\alpha\}=x$.


Семейство $\Re$ удовлетворяет следующему свойству. Пусть имеется некоторый


произвольный конусный отрезок $\langle u;v\rangle$ и каждому


$\alpha\in\Omega$ сопоставлен $x^\alpha\in\langle u;v\rangle$. Если


$\sum\limits_\alpha\{R_\alpha x^\alpha\}\in B$,


то необходимо $\sum\limits_\alpha\{R_\alpha x^\alpha\}\in


\langle u;v\rangle$.





Оператор $T$ называется гетерогенным, если для любых $u,v\in B$


\begin{equation}


\sum_\alpha\{R_\alpha T(P_\alpha u+Q_\alpha v)\}\in B


\end{equation}


и для любых $x'\ge x$, $y\ge y'$


\begin{equation}


\sum_\alpha\{R_\alpha T(P_\alpha x'+Q_\alpha y')\}\ge


\sum_\alpha\{R_\alpha T(P_\alpha x+Q_\alpha y)\}.


\end{equation}


Если $T$ оставляет инвариантным некоторое множество $M$ и условия (1), (2),


равно


как и условия, налагаемые на $\Im$, $\Re$, справедливы при дополнительном


предположении,


что все элементы, участвующие в их формулировке, принадлежат $M$,


то $T$ называется


гетерогенным на $M$.





Если $\Re$ состоит из одного тождественного преобразования


$I$ ($\Omega$ содержит лишь


одну точку $\alpha$), то условие (1) выполняется автоматически, а (2)


переходит в


\begin{equation}


T(P_\alpha x'+Q_\alpha y')\ge T(P_\alpha x+Q_\alpha y).


\end{equation}


Если при этом $P_\alpha=I$, то $T$ --- изотонный, если же $Q_\alpha=I$,


то $T$ --- антитонный оператор.





Эти определения взяты из работы \cite{Op1}.





Будем говорить далее, что семейства $\Im$ и $\Re$ обладают свойством


замкнутости, если выполняются условия


\begin{itemize}


\item[1)] из сходимости $x_n\to x$, $y_n\to y$,


$P_\alpha x_n+Q_\alpha y_n\to P_\alpha x_0+Q_\alpha y_0$


при $n\to\infty$ для


любого $\alpha\in \Omega$ следует, что $x,y\in B$ и


$P_\alpha x+Q_\alpha y=P_\alpha x_0+Q_\alpha y_0$;


\item[2)] пусть при любом $\alpha\in\Omega$ и $n\to\infty$\;


$(x^\alpha)_n\to x^\alpha$,


$\sum\limits_\alpha\{R_\alpha(x^\alpha_n)\}\in B$ при любом


$n\ge 0$,


$\sum\limits_\alpha\{R_\alpha(x^\alpha_n)\}\to


\sum\limits_\alpha\{R_\alpha(x^\alpha_0)\}$,


тогда $x^\alpha\in B$ и


$\sum\limits_\alpha\{R_\alpha(x^\alpha_0)\}=


\sum\limits_\alpha\{R_\alpha(x^\alpha)\}$.


\end{itemize}





Рассмотрим нелинейное операторное уравнение


\begin{equation}


Lx=T(x),


\end{equation}


заданное на некотором выпуклом множестве $M\subset B$


с полуупорядоченным конусом $K$.


Предположим, что оператор $T$ дифференцируем на $M$


по Фреше, и его производная $T'$


обладает свойством гетерогенности на


$\langle u_0;v_0\rangle$, где $u_0,v_0\in M$,


$\langle u_0;v_0\rangle\subset M$, --- некоторые


начальные приближения к решению уравнения (4), последующие приближения к


которому определяются алгоритмом


\begin{align}


Lu_{n+1}&=T(u_n)+\sum_\alpha\{R_\alpha T'(P_\alpha u_n+


Q_\alpha v_n)\}(u_{n+1}-u_n), \\


Lv_{n+1}&=T(v_n)+\sum_\alpha\{R_\alpha T'(P_\alpha u_n+


Q_\alpha v_n)\}(v_{n+1}-v_n).


\end{align}


Предполагается, что для любого $n$ существуют единственные решения


уравнений (5),


(6) и для расщепления $T'$ справедлива формула Ньютона-Лейбница.





\begin{thm}


Пусть для начальных приближений $u_0,v_0\in M$


выполняются неравенства


\begin{equation}


u_0\le v_0,\quad Lu_0\le T(u_0),\quad T(v_0)\le Lv_0.


\end{equation}


На $\langle u_0;v_0\rangle\subset M$ существует положительный оператор


\begin{equation}


A^{-1}=\biggl(L-\sum_\alpha\{R_\alpha T'(P_\alpha+Q_\alpha)\}


\biggr)^{-1}\ge 0.


\end{equation}





Тогда последовательные приближения, определяемые алгоритмом


{\em (5)--(6)},


удовлетворяют неравенствам


\begin{equation}


u_n\le u_{n+1}\le v_{n+1}\le v_n,\quad


n=0,1,2,\dotsc


\end{equation}


\end{thm}





{\bf Доказательство.} Вычитая $Lu_0\le T(u_0)$ из (5) при $n=0$


и учитывая линейность


оператора $A$, находим


$\Bigl(L-\sum\limits_\alpha\{R_\alpha T'(P_\alpha u_0+Q_\alpha v_0)\}


\Bigr)(u_1-u_0)\ge 0$.


Согласно (8)


\begin{equation}


u_0\le u_1.


\end{equation}


Аналогично доказывается соотношение


\begin{equation}


v_1\le v_0.


\end{equation}





Докажем, что $u_1\le v_1$. Вычитая из (6) соотношение (5) при


$n=0$, имеем


$$


L(v_1-u_1)=T(v_0)-T(u_0)+\sum_\alpha


\{R_\alpha T'(P_\alpha u_0+Q_\alpha v_0)\}(v_1-v_0)-


\sum_\alpha\{R_\alpha T'(P_\alpha u_0+Q_\alpha v_0)\}


(u_1-u_0).


$$


Применяя формулу Ньютона-Лейбница к разности $T(v_0)-T(u_0)$, получим


\begin{gather*}


\biggl(L-\sum_\alpha\{R_\alpha T'(P_\alpha u_0+Q_\alpha v_0)\}\biggr)


(v_1-u_1)= \\


=\int^1_0\biggl[\sum_\alpha\{R_\alpha T'(P_\alpha(u_0+\tau


(v_0-u_0))+Q_\alpha(u_0+\tau(v_0-u_0)))\}-


\sum_\alpha\{R_\alpha T'(P_\alpha u_0+Q_\alpha v_0)\}\biggr]\times \\


\times (v_0-u_0)d\tau.


\end{gather*}


Согласно первому неравенству в (7) из гетерогенности $T'$ и (8) имеем


\begin{equation}


u_1\le v_1.


\end{equation}


Неравенства (10)--(12) доказывают заключение теоремы при $n=0$.





Покажем, что элементы $u_1,v_1\in\langle u_0;v_0\rangle$


удовлетворяют второму и третьему


неравенствам в (7). Действительно,


\begin{gather*}


Lu_1-T(u_1)=T(u_0)-T(u_1)+


\sum_\alpha\{R_\alpha T'(P_\alpha u_0+Q_\alpha v_0)\}(u_1-u_0)= \\


=-\int^1_0 T'(u_0+\tau(u_1-u_0))(u_1-u_0)d\tau+


\sum_\alpha\{R_\alpha T'(P_\alpha u_0+Q_\alpha v_0)\}(u_1-u_0)=\\


=\int^1_0\biggl[\sum_\alpha\{R_\alpha T'(P_\alpha u_0+Q_\alpha v_0)\}-


\sum_\alpha\{R_\alpha T'(P_\alpha(u_0+\tau(u_1-u_0))+


Q_\alpha(u_0+\tau(u_1-u_0)))\}\biggr]\times \\


\times (u_1-u_0)d\tau.


\end{gather*}


Учитывая соотношения (10) и гетерогенность $T'$, получаем


$Lu_1\le T(u_1)$. Аналогично


доказывается третье неравенство в (7). Включение


$\langle u_1;v_1\rangle\subset\langle u_0;v_0\rangle$


гарантирует


выполнение соотношения (8). Таким образом, все условия теоремы


выполнены для элементов $u_1$, $v_1$.





Доказательство теоремы завершается индукцией по $n$. $\quad\square$





\begin{thm} Пусть $K$ --- правильный конус, операторы $T$ и $A$ замкнуты на


$\langle u_0;v_0\rangle$ и


выполняются условия теоремы {\em 1.}





Тогда уравнение {\em (4)} имеет на отрезке


$\langle u_0;v_0\rangle$ по крайней мере одно решение.


Среди всех решений имеется наименьшее $u^*$ и наибольшее $v^*$,


к которым сходятся


последовательные приближения, определяемые алгоритмом


{\em (5)--(6)}, причем


\begin{equation}


u_n\le u^*\le v^*\le v_n,\quad n=0,1,2,\dotsc


\end{equation}


\end{thm}





\begin{pf} Так как по теореме 1 последовательности $\{u_n\}$ и $\{v_n\}$


монотонны


и ограничены, а $K$ --- правильный конус, то существуют


$\lim\limits_{n\to\infty}u_n=u^*$,


$\lim\limits_{n\to\infty}v_n=v^*$, и согласно


(9) $u^*\le v^*$. Переходя к пределу в (5)--(6),


учитывая замкнутость операторов


$A$ и $T$, убеждаемся, что $u^*$ и $v^*$


являются решениями уравнения (4). Если $x^*$ ---


произвольное решение этого уравнения на


$\langle u_0;v_0\rangle$, то элементы $u_0$, $x^*$ удовлетворяют всем


условиям теоремы, и по доказанному $u_n\le x^*$.


Переходя к пределу в этом неравенстве,


убеждаемся, что $u^*\le x^*$, т.е. $u^*$ --- наименьшее решение.


Аналогично доказывается, что


$v^*$ --- наибольшее решение.





Докажем последнее утверждение теоремы, проведя один шаг индукции. Пусть


неравенства (13) имеют место для некоторого $n=m>0$.


Так как $u^*$ --- решение уравнения


(4), то из (5) получаем


$$


Lu_{m+1}-Lu^*=T(u_m)-T(u^*)+\sum_\alpha\{R_\alpha T'


(P_\alpha u_m+Q_\alpha v_m)\}(u_{m+1}-u_m),


$$


или


\begin{gather*}


\biggl(L-\sum_\alpha\{R_\alpha T'(P_\alpha u_m+Q_\alpha v_m)\}\biggr)


(u_{m+1}-u^*)= \\


\int^1_0\biggl[\sum_\alpha\{R_\alpha T'(P_\alpha(u^*+\tau(u_m-u^*))+


Q_\alpha(u^*+\tau(u_m-u^*)))\}-


\sum_\alpha\{R_\alpha T'(P_\alpha u_m+Q_\alpha v_m)\}\biggr]\times \\


\times (u_m-u^*)d\tau\le 0.


\end{gather*}


Следовательно, $u_{m+1}\le u^*$. Аналогично доказывается,


что $v^*\le v_{m+1}$. Таким образом,


неравенства (13) доказаны для любого $n$.


\end{pf}





\begin{thm} Пусть выполняются условия теоремы {\em 1.}


Конус $K$ нормален. Оператор\linebreak


$T-\sum\limits_\alpha\{R_\alpha T'(P_\alpha+Q_\alpha)\}$


вполне непрерывен, $A^{-1}$ непрерывен на $\langle u_0;v_0\rangle$.





Тогда уравнение {\em (4)} имеет на


$\langle u_0;v_0\rangle$ по крайней мере одно решение $x^*$, к


которому сходятся последовательности $\{u_n\}$, $\{v_n\}$, определяемые


алгоритмом {\em (5)--(6)},


причем


\begin{equation}


u_n\le u_{n+1}\le x^*\le v_{n+1}\le v_n,\quad n=0,1,2,\dotsc


\end{equation}


\end{thm}





\begin{pf} Монотонность последовательных приближений доказана в


теореме 1. Для доказательства существования решения


$x^*\in\langle u_n;v_n\rangle$, $n=0,1,2,\dotsc$,


рассмотрим оператор


\begin{equation}


\Gamma x=x+A^{-1}(T(x)-Lx)\equiv A^{-1}


\biggl[T(x)-\sum_\alpha\{R_\alpha T'(x)\}\biggr].


\end{equation}


Покажем, что $\Gamma$ --- изотонный оператор на


$\langle u_0;v_0\rangle$. Действительно, если


$x,y\in\langle u_0;v_0\rangle$ и


$x\le y$, то


\begin{gather*}


\Gamma x-\Gamma y=x-y+A^{-1}(T(x)-T(y)-L(x-y))=\\


=x-y+A^{-1}\biggl[\int^1_0 T'(y+\tau(x-y))(x-y)d\tau-L(x-y)\biggr]=\\


=x-y+A^{-1}\biggl[ \int^1_0\biggl(\sum_\alpha\{R_\alpha T'


(P_\alpha(y+\tau(x-y))+Q_\alpha(y+\tau(x-y)))\}\biggr)


(x-y)d\tau-L(x-y)\biggr].


\end{gather*}


Так как $u_0\le x\le y\le v_0$, то


$$


\sum_\alpha\{R_\alpha T'(P_\alpha(y+\tau(x-y))+Q_\alpha(y+\tau(x-y)))\}


(x-y)\le\sum_\alpha\{R_\alpha T'(P_\alpha u_0+Q_\alpha v_0)\}


(x-y)


$$


и


$$


\Gamma x-\Gamma y\le x-y-A^{-1}\biggl[L-\sum_\alpha


\{R_\alpha T'(P_\alpha u_0+Q_\alpha v_0)\}(x-y)\biggr]=0,


$$


т.\,е. $\Gamma x\le \Gamma y$.


Из полной непрерывности


$T-\sum\limits_\alpha\{R_\alpha T'(P_\alpha+Q_\alpha)\}$


и непрерывности $A^{-1}$ следует полная


непрерывность оператора $\Gamma$. Так как


$$


\Gamma u_n=u_n+A^{-1}(T(u_n)-L(u_n))\ge u_n,\quad


\Gamma v_n=v_n+A^{-1}(T(v_n)-L(v_n))\le v_n,\quad


n=0,1,2,\dotsc,


$$


то оператор $\Gamma$ отображает отрезок


$\langle u_n;v_n\rangle$ в себя и по теореме Шаудера имеет на


этом отрезке неподвижную точку $x^*$, для которой $Lx^*=T(x^*)$.


Неравенство (14)


следует из принадлежности $x^*$ отрезку


$\langle u_n;v_n\rangle$, $n=0,1,2,\dotsc,$ и монотонности


$\{u_n\}$, $\{v_n\}$.


\end{pf}





\begin{thm} Пусть выполняются условия теоремы {\em 1.}


Если справедливо хотя бы


одно из следующих условий:


\begin{itemize}


\item[а)] конус $K$ сильно миниэдральный;


\item[б)] конус $K$ телесный, нормальный и миниэдральный, оператор


$A^{-1}\Bigl[T-\sum\limits_\alpha\{R_\alpha T'


(P_\alpha+Q_\alpha)\}\Bigr]$ компактный на $\langle u_0;v_0\rangle$,


\end{itemize}


то на $\langle u_0;v_0\rangle$ существует по крайней мере одно решение


$x^*$ уравнения {\em (4)}, удовлетворяющее условиям {\em (14).}


\end{thm}





\begin{pf} Покажем, что оператор $\Gamma$ в (15) имеет неподвижную точку


$x^*\in\langle u_0;v_0\rangle$.


Тогда из (15) будет следовать, что $x^*$ --- решение уравнения (4).





Пусть выполняется случай а). Обозначим через $\Re$ множество таких


элементов $x\in\langle u_0;v_0\rangle$,


что $\Gamma x\ge x$. Это множество не пусто, т.к.


$\Gamma u_0=u_0+A^{-1}(T(u_0)-Lu_0)\ge u_0$.


Поскольку $\Gamma(\Gamma x)\ge\Gamma x$, $x\in\Re$, в силу изотонности


$\Gamma$, то


$\Gamma\Re\subset\Re$. Пусть $x^*=\sup\Re$. Очевидно,


$u_0\le x^*\le v_0$, т.\,е.


$x^*\in\langle u_0;v_0\rangle$. Из изотонности $\Gamma$


и определения $x^*$ вытекает, что


$\Gamma x^*\ge\Gamma x\ge x$ для любого $x\in\Re$. Это значит, что


$\Gamma x^*$ является верхней границей для $\Re$.


Поэтому $\Gamma x^*\ge x^*$. Но тогда $x^*\in\Re$, в силу чего


$\Gamma x^*\in \Re$ и $\Gamma x^*\le x^*$. Итак, $\Gamma x^*=x^*$,


откуда $Lx^*=T(x^*)$.


Соотношения (14) следуют из


инвариантности отрезков


$\langle u_n;v_n\rangle$, $n=0,1,2,\dotsc,$ для оператора $\Gamma$


и монотонности $\{u_n\}$, $\{v_n\}$.





Пусть выполняется случай б). Множество


$\Re=\Gamma\langle u_0;v_0\rangle$ компактное


и $u_1=\Gamma u_0\in\Re$. Обозначим через $\Re_1$


множество тех элементов из $\Re$, для которых


$\Gamma x\ge x$. Очевидно, $u_1\in\Re$, так что $\Re_1$ не пусто.


Так как это множество компактно и


ограничено сверху, то существует $x^*=\sup\Re_1$ [2].


Тогда согласно доказательству


случая а) $Lx^*=T(x^*)$.


\end{pf}





Условия а), б) теоремы 4 существенно различны. Действительно, например,


конус


неотрицательных функций в $C[0;1]$


является телесным, нормальным и миниэдральным,


не являясь в то же время сильно миниэдральным [1].





\begin{thm} Пусть кроме условий теоремы {\em 2} выполняются условия


\begin{itemize}


\item[1)] конус $K$ нормален$;$


\item[2)] оператор $A^{-1}$


ограничен на $\langle u_0;v_0\rangle$, $\|A^{-1}\|\le B;$


\item[3)] оператор


$\sum\limits_\alpha\{R_\alpha T'(P_\alpha+Q_\alpha)\}$


удовлетворяет условию Липшица по обеим переменным


\begin{align*}


\biggl\|\sum_\alpha\{R_\alpha T'(P_\alpha u+Q_\alpha z)\}-


\sum_\alpha\{R_\alpha T'(P_\alpha v+Q_\alpha z)\}\biggr\|


&\le L_1\|u-v\|,\\


\biggl\|\sum_\alpha\{R_\alpha T'(P_\alpha z+Q_\alpha u)\}-


\sum_\alpha\{R_\alpha T'(P_\alpha z+Q_\alpha v)\}\biggr\|


&\le L_2\|u-v\|;


\end{align*}


\item[4)] $B(L_1+L_2)\|u_0-v_0\|<2$.


\end{itemize}





Тогда последовательности $\{u_n\}$, $\{v_n\}$


сходятся соответственно снизу и сверху к


единственному на $\langle u_0;v_0\rangle$


решению $x^*=u^*=v^*$ уравнения {\em (4)}, и скорость их сходимости


характеризуется неравенствами


\begin{equation}


\begin{split}


\|x^*-u_n\|&\le I^{2^n-1}\|v_0-u_0\|^{2^n},\\


\|v_n-x^*\|&\le I^{2^n-1}\|v_0-u_0\|^{2^n},


\end{split}


\end{equation}


где $I=\frac{1}{2}B(L_1+L_2)$.


\end{thm}





\begin{pf} Вычитая (5) из (6), приходим к равенству


\begin{gather*}


v_{n+1}-u_{n+1}=


A^{-1}\int^1_0\biggl[\sum_\alpha\{R_\alpha


T'(P_\alpha(u_n+\tau(v_n-u_n))+


Q_\alpha(u_n+\tau(v_n-u_n)))\}- \\


-\sum_\alpha\{R_\alpha T'(P_\alpha u_n+Q_\alpha v_n)\}\biggr]


(v_n-u_n)d\tau.


\end{gather*}


Оценивая его по норме в силу условий 2)--4) теоремы, получим


\begin{multline}


\|v_{n+1}-u_{n+1}\|\le\|A^{-1}\|\int^1_0


(L_1\tau\|v_n-u_n\|+L_2(1-\tau)\|v_n-u_n\|)


\|v_n-u_n\|d\tau\le \\


\le B\frac{L_1+L_2}{2}\|v_n-u_n\|^2=I\|v_n-u_n\|^2.


\end{multline}


Положим в (17) последовательно $n=0,1,2,\dotsc$ Тогда


$$


\|v_1-u_1\|\le I\|v_0-u_0\|^2,\quad


\|v_2-u_2\|\le I\|v_1-u_1\|^2\le


I^3\|v_0-u_0\|^4.


$$


По индукции заключаем, что


$\|v_n-u_n\|\le I^{2^n-1}\|v_0-u_0\|^{2^n}$. Согласно условию 4) теоремы и


заключению теоремы 2\;


$\lim\limits_{n\to\infty}v_n=\lim\limits_{n\to\infty}u_n=x^*$ ---


единственное


решение уравнения (4). Из


нормальности конуса и из неравенств


$0\le x^*-u_n\le v_n-u_n$, $0\le v_n-x^*\le v_n-u_n$, которые


непосредственно вытекают из заключения теоремы 2, получаем оценки (16).


\end{pf}





Использование алгоритмов (5)--(6) требует решения линейных операторных


уравнений на каждом шаге вычислительного процесса, в то время как


построение


последовательных приближений по соответствующим модифицированным


алгоритмам


\begin{align*}


Lu_{n+1}&=T(u_n)+\sum_\alpha\{R_\alpha T'(P_\alpha u_0+


Q_\alpha v_0)\}(u_{n+1}-u_n),\\


Lv_{n+1}&=T(v_n)+\sum_\alpha


\{R_\alpha T'(P_\alpha u_0+Q_\alpha v_0)\}


(v_{n+1}-v_n),\quad n=0,1,2,\dotsc,


\end{align*}


предполагает обращение линейного оператора


$L-\sum\limits_\alpha\{R_\alpha T'(P_\alpha u_0+Q_\alpha v_0)\}$


лишь в одной


точке $\langle u_0;v_0\rangle$.


При этом квадратичная скорость сходимости понижается до скорости


сходимости геометрической прогрессии.





В качестве примера рассмотрим в $R^m$ систему нелинейных алгебраических


уравнений


\begin{equation}


x_k=f_k(x_1,x_2,\dotsc,x_m),\quad k=1,\dotsc,m,


\end{equation}


где функции $f_k(x)$ дифференцируемые и каждая частная производная


$\dfrac{\partial f_i(x)}{\partial x_j}$, $i,j=1,\dotsc,m$,


по одним переменным монотонно возрастает (не убывает), а по остальным


переменным монотонно убывает (не возрастает). С каждой частной производной


$\dfrac{\partial f_i(x)}{\partial x_j}$ свяжем два подмножества индексов


$G_{ij}\subset I$ и $H_{ij}\subset I$


($G_{ij}\cup H_{ij}=I$, $G_{ij}\cap H_{ij}=0$), $I=\{i\mid


j=1,\dotsc,m\}$.


Номер $k\in G_{ij}$, если


$\dfrac{\partial f_i(x)}{\partial x_j}$ не убывает по $x_k$, и


$k\in H_{ij}$ в противном случае.





Пусть $P_{ij}$ --- проектор, задаваемый матрицей


$P_{ij}=(p^{ij}_{\mu\nu})$, где $p^{ij}_{\mu\nu}=1$, если $\mu\in G_{ij}$,


остальные $p^{ij}_{\mu\nu}=0$. Проектор $Q_{ij}$ определяется матрицей


$E-P_{ij}$, где $E$ --- единичная


матрица.





Пусть оператор $R_i$ из семейства $R=\{R_i\mid i\in I\}$


сопоставляет вектору $x\in R^m$ его


$i$-ю координату $x_i$, т.е. $R_ix=x_i$,


а операция $\sum\limits_i$ сопоставляет множеству компонент


$\{x_i\}$ соответствующий вектор $x=\sum\limits_i\{R_ix\}$.





Введем в $R^m$ полуупорядоченность с помощью неотрицательного ортанта


$R^m_+$, тогда


оператор


$\dfrac{\partial f}{\partial x}=\Bigl(


\dfrac{\partial f_i}{\partial x_j}\Bigr)^m_{i,j=1}$


будет удовлетворять условиям гетерогенности, и система


(5)--(6) примет вид


\begin{equation}


\begin{split}


u_{n+1}&=f(u_n)+\sum_i\biggl\{


\sum_j(R_j(R_if)')(P_{ij}u_n+Q_{ij}v_n)\biggr\}(u_{n+1}-u_n),\\


v_{n+1}&=f(v_n)+\sum_i\biggl\{\sum_j


(R_j(R_if)')(P_{ij}u_n+Q_{ij}v_n)\biggr\}(v_{n+1}-v_n).


\end{split}


\end{equation}


Обозначим


$A(u_n;v_n)=\sum\limits_i\Bigl\{\sum\limits_j


(R_j(R_if)')(P_{ij}u_n+Q_{ij}v_n)\Bigr\}$.





\begin{thm} Пусть существуют $u_0,v_0\in R^m$ такие, что


$$


u_0\le f(u_0),\quad f(v_0)\le v_0,\quad u_0\le v_0


$$


и на $[u_0;v_0]$ выполняются условия


\begin{enumerate}


\item[1)] матрица $A$ гетерогенная{\em ;}


\item[2)] для элементов $b_{jk}$, $j,k=1,\dotsc,m$,


матрицы $B=\Delta-A$, где $\Delta$ --- матрица


Кронекера, выполнено хотя бы одно из условий


\begin{enumerate}


\item[а)] $b_{jk}\le 0$ для $j\ne k$, $B$ неразложима,


существуют $y>0$ и $r\ge 0$, $r\not\equiv 0$, такие, что $By=r$;


\item[б)] $b_{jj}\ge 0$, $b_{jk}\le 0$ при $j\ne k$,


выполнен слабый признак сумм по строкам, и матрица


$B$ неразложима, либо выполнен сильный признак сумм по строкам;


\item[в)] $a_{jk}\ge 0$, $j,k=1,\dotsc,m$, где $a_{jk}$ ---


элементы матрицы $A$, $\|A\|<1$;


\item[г)] $B$ --- симметричная, положительно определенная


матрица и $a_{jk}\ge 0$ для любых $j$,~$k$.


\end{enumerate}


\end{enumerate}





Тогда система {\em (18)} на $[u_0;v_0]$


имеет по крайней мере одно решение $x^*$, к


которому сходятся последовательности $\{u_n\}$, $\{v_n\}$,


определяемые {\em (19),} и $u_n\le u_{n+1}\le x^*\le v_{n+1}\le v_n$.


\end{thm}





Каждое из условий а), б), в), г) п.\,2 теоремы гарантирует выполнение


условия (8) теоремы~1 ([7], \S\,23).





\begin{thebibliography}{9}





\bibitem{Kra}


Красносельский М.А. {\em Положительные решения операторных уравнений}.


-- М.: Физматгиз, 1962. -- 394~с.





\bibitem{Ste}


Стеценко В.Я. {\em О неподвижных точках нелинейных отображений} //


Сиб. матем. журн. -- 1969. -- Т.\,10. -- \No\,3. -- С.\,642--652.





\bibitem{Ba}


Бахтин И.А. {\em О существовании общих неподвижных точек для абелевых


совокупностей разрывных операторов} // Сиб. матем. журн. -- 1972. --


Т.\,13. -- \No\,2. -- С.\,243--251.


\bibitem{Op1}


Опойцев В.И. {\em Гетерогенные и комбинированно-вогнутые операторы}


// Сиб. матем. журн. -- 1975. -- Т.\,16. -- \No\,4. -- С.\,781--792.


\bibitem{Op2}


Опойцев В.И. {\em Обобщение теории монотонных и вогнутых операторов} //


Тр. Моск. матем. о-ва. -- 1978. -- Т.\,36. -- С.\,237--273.


\bibitem{Kur}


Курпель Н.С., Шувар Б.А. {\em Двусторонние операторные неравенства и их


применения}. -- Киев: Наук. думка, 1980. - 267~с.


\bibitem{Kol}


Коллатц Л. {\em Функциональный анализ и вычислительная математика}.


-- М.: Мир, 1969. -- 447~с.





\end{thebibliography}





\vskip 0.5cm





\post{\it Черниговский педагогический}{\it Поступила}


\post{\it институт $($Украина$)$}{19.01.1995}





\label{end}


\end{document}


